\section{Introduction}
\label{sec:introduction}

Tropical peat swamp forests are among the most carbon-dense ecosystems on the planet, sequestering up to twenty times more organic carbon per unit area than adjacent upland tropical rainforests on mineral soils \citep{Page2002Nature, Dommain2011ProcRSocB}. In undisturbed equatorial environments, high phreatic water tables remain near or above the surface throughout the annual cycle, creating anoxic conditions that retard organic matter decomposition and render natural peat formations virtually non-combustible \citep{Wosten2008Geoderma, Dadap2021NatComm}. Over the past three decades, however, Southeast Asian peatlands---predominantly concentrated across Sumatra and Borneo---have experienced severe anthropogenic disruption driven by industrial forestry, oil palm expansion, and large-scale state-sponsored canal excavation schemes \citep{Miettinen2016ERL, Dohong2017RestorEcol, Langner2007IJWF, Carlson2012PNAS}.

The most catastrophic manifestation of this ecological transformation occurred in Central Kalimantan (\textit{Kalimantan Tengah}), Indonesia, through the execution of the Mega Rice Project (\textit{Proyek Lahan Gambut} / ex-PLG) between 1995 and 1999 \citep{Silvius2005Wetlands}. Over 4,000\,km of primary, secondary, and tertiary drainage canals were excavated across more than one million hectares of ombrotrophic peat domes between the Sebangau, Kahayan, Kapuas, and Barito river systems \citep{Ritzema2014HESS}. This artificial drainage infrastructure fundamentally dismantled the hydrological equilibrium of the landscape: canals act as continuous conduits that accelerate gravitational discharge during dry seasons, drawing the groundwater table well below the legally mandated threshold of $-40$\,cm established under Indonesian Government Regulations (PP No.~71/2014 and PP No.~57/2016). When meteorological dry spells coincide with positive El Ni\~no--Southern Oscillation (ENSO) anomalies or positive Indian Ocean Dipole (pIOD) events, the desiccated surface peat transitions into an acutely combustible fuel matrix \citep{Cochrane2003Nature, Hooijer2012Biogeosciences, Field2016PNAS}. Once ignited, low-temperature subterranean smoldering combustion propagates horizontally and vertically through the catotelm layer for months, releasing ancient carbon sequestered over millennia \citep{Usup2004Tropics, Wooster2018NatGeosci, Vetrita2020RemoteSensing}.

The atmospheric and socio-economic ramifications of recurrent Indonesian peatland wildfires are devastating on regional and global scales. Catastrophic fire episodes in 1997, 2015, and 2019 released an estimated 0.81--2.57\,Gt, 1.2\,Gt, and 0.71\,Gt of $\text{CO}_2$-equivalent emissions, respectively, momentarily exceeding the daily fossil fuel emissions of major industrialized nations \citep{Page2002Nature, Field2016PNAS, VanderWerf2017ESSD}. Simultaneously, peat smoldering generates massive quantities of toxic fine particulate matter ($\text{PM}_{2.5}$), carbon monoxide, and aerosol precursors that induce transboundary smoke-haze crises across Indonesia, Singapore, and Malaysia, culminating in acute public health emergencies and substantial regional economic losses \citep{Gaveau2014PLOSONE, Kiely2021ACP, Graham2022SciRep}.

Despite extensive scientific documentation, operational wildfire risk prediction and policy interventions in tropical peatlands remain fundamentally impeded by two major dilemmas:

First, \textbf{the methodological pitfall of spatial autocorrelation and spatial memorization in machine learning}. The remote sensing and geospatial hazard communities have increasingly adopted supervised machine learning algorithms---such as Random Forests, Gradient Boosted Trees, and Deep Neural Networks---for active fire mapping and susceptibility prediction \citep{Breiman2001ML}. However, environmental and climatic processes operate under Tobler's First Law of Geography: ``everything is related to everything else, but near things are more related than distant things'' \citep{Tobler1970EconGeogr, Legendre1993Ecology}. When spatial hazard models incorporate geodetic coordinates (\texttt{latitude} and \texttt{longitude}) or administrative names as explicit predictor variables, decision trees frequently construct decision boundaries that partition geographic space directly \citep{Meyer2018EcolModel, Dormann2007Ecography}. When evaluated using conventional random cross-validation ($k$-fold or random holdout train-test splits), spatial proximity between training and testing observations leads to pervasive spatial data leakage \citep{Roberts2017Ecography, Wadoux2021EcolModel}. The algorithm essentially memorizes local spatial clusters rather than learning genuine biophysical mechanisms governing combustion. When deployed operationally to predict fire behavior in unfamiliar administrative units or across out-of-sample regions, the predictive skill of such models collapses catastrophically \citep{Meyer2018EcolModel, Schratz2019EcolModel, MeyerPebesma2021NatComm}.

Second, \textbf{the land tenure governance void and polarized culpability narratives}. Public, legal, and non-governmental discourse surrounding Indonesian peat fires frequently oscillates between two contradictory positions: either attributing the majority of ignitions to intentional corporate arson by palm oil and pulpwood concessionaires \citep{Tacconi2003CIFOR, Carlson2012PNAS}, or conversely framing wildfires entirely as the outcome of unregulated traditional slash-and-burn smallholder practices and illegal encroachment \citep{Cattau2016GCB}. Rigorous, spatially explicit matching of high-resolution satellite thermal detections against verified, cadastrally audited corporate concession perimeters remains sparse in the peer-reviewed literature, obscuring the empirical reality of land tenure governance across degraded tropical peatlands \citep{Putranto2023LandUsePol}.

To resolve these pressing methodological and governance challenges, this study conducts an empirical investigation of the wildfire crisis in Central Kalimantan. Drawing on 43,937 satellite-detected active fire hotspots from the NASA FIRMS constellation during the October 2026 fire pulse, we test two explicit research hypotheses:
\begin{enumerate}
    \item[$\mathbf{H_1}$:] \textit{Inclusion of geodetic coordinates in tree-based fire prediction models produces an optimistic illusion of predictive skill under standard random cross-validation due to spatial data leakage, but results in near-complete performance collapse under spatially disjoint block cross-validation due to spatial memorization.}
    \item[$\mathbf{H_2}$:] \textit{A pure biophysical feature representation---grounded strictly in hydrological canal proximity, edaphic peat depth, antecedent precipitation deficits, and fire weather indices---retains generalizable predictive power and positive precision across unseen geographic regions.}
\end{enumerate}

To evaluate these hypotheses and provide actionable policy insights, we present: (1) a decadal analysis of official Ministry of Environment and Forestry (SiPongi KLHK) burned area records (2018--2026) alongside micro-scale diurnal satellite fire dynamics; (2) an exhaustive spatial join against 627 certified palm oil concessions (HGU) and 185 commercial forestry licenses (PBPH) to quantify land tenure distributions; (3) a comparative machine learning benchmark contrasting a coordinate-dependent Spatial Memorization Model (Model~A) against a coordinate-free Pure Biophysical Model (Model~B) under 5-fold Spatial Block Cross-Validation (\texttt{GroupKFold} by regency); and (4) an interpretability analysis using exact TreeSHAP algorithms to discover critical non-linear biophysical tipping points governing severe Fire Radiative Power ($\text{FRP} \ge 12.0$\,MW).
